\paragraph{Internal validity.}
The measures are computed by one engine and the result is a function of the engine's reading of the supervisor's text. The reading was checked against the text of Annex 1 for the sizes, the floors, the pair rule, the charge, the grid and the thresholds, and the contract tests of the engine reproduce its golden values; there is no second implementation of Annex 1 to reconcile with, so a defect common to the text's reading and the engine would be invisible. The reading of the floor provision that keeps a pre-shock rate below the floor is settled by the plural of the text (``estos pisos'') but is still a reading. The comparison across seeds measures the variation of the synthetic balance sheet; it does not measure the uncertainty of the curve, which is one fit of one day, with four points per curve and a root-mean-square error of \rmseClp{} and \rmseUr{} basis points, and with a policy rate in place of a missing peso yield. The yields are read as semiannual bond-equivalent par yields, which has not been checked against the central bank's methodology note.

\paragraph{External validity.}
Banco Andino is synthetic. Its totals are calibrated to the banking system, but the UF share of each portfolio, the repricing profile of each item, the deposit split and every behavioural parameter are assumptions, so the proportions of the Shapley values describe the treatment on this book and not any bank. The cancellation of the own curve and the own sizes in the value measure depends on where the UF flows lie in the grid; a bank with other maturities would have other weights. The net UF positions at which the two treatments classify more than half of the books differently, \sweepMajLo{} to \sweepMajHi{} percent of level 1 capital, are at or above the whole capital of the bank, and \sweepGridHi{} percent is the top of the range tested, so the sweep does not show where the difference ends; single books differ from \sweepFirstDiff{} percent. It is not known whether any Chilean bank holds a position of that size, which the aggregates published by the supervisor do not state. The curves are those of one day, September 2026; periods of negative real rates, in which the floor of minus two percent was designed to matter, are not in the raw files, and the floor experiment replaces them by a translation of the fitted curve, which is a what-if and not a history.

\paragraph{Construct validity.}
The paper measures a rule, not an exposure. The Annex 1 measure is a standardised supervisory number; it assumes a run-off balance for the value measure and a constant balance for the income measure, ignores the optionality that the text treats separately, and applies one set of behavioural factors. A difference between the Annex 1 and the single-currency measures is not a difference in the risk of a bank. The single-currency treatment is a counterfactual that a bank aware of indexed units would not use, and the comparison shows what the separate treatment adds to a number; it does not show that a bank that used the counterfactual would misjudge its risk. The indexation charge is read as a stress of 2.5 percentage points in either direction, and the paper does not model the process of inflation or the one-month lag of the UF accrual. The attribution by the Shapley value depends on the choice of five switches; a different partition of the treatment, the curve and the sizes as one switch for example, would give other values, and the paper does not claim that the five are the only decomposition.

\paragraph{Scope.}
The measurement is for a bank whose significant currencies are the peso and the UR, and the aggregation across further currencies is not studied. The option treatment of the text is not implemented. The hedges and the dynamic balance sheet are outside the standardised measures. The statements about the absence of a public implementation of Annex 1 and of a published decomposition are limited to the search behind this work.
